\section{Conclusion}
\label{sec:conclusion}

% draft -- finalise later

This survey has reviewed the emerging field of deep learning-based reconstruction for quantum imaging systems, covering ghost imaging, compressed sensing with quantum correlations, sub-shot-noise imaging, and hybrid classical--quantum architectures.

Our principal findings are as follows. First, deep learning enables measurement reductions of one to two orders of magnitude in ghost imaging reconstruction, with end-to-end learned approaches providing the largest gains through joint optimisation of measurement and reconstruction. Second, the field is converging towards U-Net and GAN-based architectures, mirroring trends in classical computational imaging. Third, critical gaps remain in standardised benchmarking, high-resolution scalability, and systematic evaluation of quantum advantage in the presence of deep learning.

% add: research gaps summary, future directions, community recommendations

The intersection of quantum imaging and deep learning remains a fertile area for research, with significant potential for both fundamental insights and practical applications. As quantum optical hardware matures and deep learning methods continue to advance, their synergistic integration promises imaging capabilities that transcend the limits of either approach alone.
